\documentclass{article}
\usepackage[spanish, activeacute]{babel}
\usepackage[utf8]{inputenc}

\usepackage{graphicx}
\usepackage{float}
\usepackage{listings}
\usepackage{xcolor}
\usepackage{ dsfont }
\usepackage{vmargin}
\usepackage{ wasysym }
\usepackage{booktabs}
\usepackage{ amssymb }
\usepackage{hyperref}
\usepackage{tikz}

\pagenumbering{roman}
\setpapersize{A4}
\setmargins{2cm}       % margen izquierdo
{1.5cm}                        % margen superior
{17cm}                      % anchura del texto
{23.42cm}                    % altura del texto
{9pt}                           % altura de los encabezados
{1cm}                           % espacio entre el texto y los encabezados
{0pt}                             % altura del pie de página
{2cm}                           % espacio entre el texto y el pie de página

\title{Práctica 4\\Primer caso de uso - Login}
\author{
    Diana Laura Luciano Paniagua\\
    \texttt{dianaluciano@ciencias.unam.mx}
    \and
    Luis Fernando Yang Fong Baeza\\
    \texttt{fernandofong@ciencias.unam.mx}
    \and
    Valeria García Landa\\
    \texttt{vale.garcia.landa@gmail.com}
    \and
    Francisco Valdés Souto\\
    \texttt{fvaldes@ciencias.unam.mx}
}
\date{\today}

\begin{document}

\maketitle

\section{Introducción}

Un \textit{login} es la base para cualquier sistema web puesto que su funcionalidad puede depender de la autentificación de un usuario, ya que esto relaciona directamente con un \textit{actor} de cualquier caso de uso. El propósito de esta práctica es que el alumno haga la base (un login) para su sistema web que entregará al final del curso.\\

Este \textit{login} tiene que ser utilizando las herramientas que se vieron durante el laboratorio, o sea, Python en versión $\geq 3.12.x$, MySQL $\geq 9.5.x$, Flask SQL Alchemy y Flask.\\

\textbf{Nota:} Aunque el sistema web se entregue de manera grupal, cada alumno debe de entregar su propia práctica para aprender los conocimientos necesarios para implementar un caso de uso tan importante y que abarca tanto como lo es un \textit{login}. Posteriormente cada equipo deberá decidir cuál es la mejor base para empezar su proyecto y construir sobre él.

\subsection{Entregados}

Los entregados para esta práctica son:

\begin{enumerate}
\item PDF de la práctica 
\item Plantilla de \LaTeX práctica para responder las preguntas.
\item Archivo de \texttt{.sql}
\end{enumerate}

\section{Desarrollo}

\subsection{Preparación de la práctica}

En esta ocasión, no se le entrega un código fuente al alumno puesto que es su responsabilidad levantar un proyecto desde cero, al igual que es su deber levantar la base de datos que les fue proporcionada por los ayudantes de laboratorio, esto para un mejor aprendizaje.

\subsection{Realización de la práctica}

Para realizar la práctica, el alumno deberá de crear una aplicación en Flask que muestre una pantalla de inicio, que capture al menos dos datos, usuario y contraseña y que dicha pantalla permita enviar la información capturada con un botón. En caso de que el usuario y contraseña sean correctos, la aplicación debe de presentar una siguiente pantalla que debe de mostrar el nombre del usuario y el rol si es que la información ingresada fue correcta, como mínimo. La validez de la información tiene que ser con respecto a lo que está definido en el script de SQL proporcionado por los ayudantes de laboratorio. El estilo de las pantallas (CSS) es a completo criterio del alumno siempre y cuando no se caiga en obscenidades.\\

En el script de la base de datos que se les fue proporcionada, se incluyen dos usuarios pero con roles diferentes, un usuario y un admin, el sistema tiene que diseñar una manera de diferenciar ambos roles, ya sea con vistas diferentes o inyectando diferentes valores pero tiene que ser una diferencia visible. En caso de que haya un error en usuario, contraseña o ambos, el sistema debe de también lograr representarlo en una vista pero permitiendo al usuario volver al inicio del flujo.\\

Por último, también es obligación del alumno solicitar a alguna herramienta de IA que le cree un HTML y CSS para sus vistas o \textit{templates} de su proyecto. Si el alumno quiere hacer a mano sus archivos HTML o CSS, es a gusto propio, esto no será evaluable.\\
 
\subsection{Verificación de la práctica}

Para verificar su práctica, el alumno deberá de actuar como un usuario final para su propio sistema y no deberá saltar ningún error inesperado, ni pantalla de error 5XX ni pantalla de \textit{Internal System Error}, deberá de cumplir correctamente con los siguientes 3 flujos:

\begin{enumerate}
    \item Inicio de sesión exitoso para usuario.
    \item Inicio de sesión exitoso para administrador.
    \item Inicio de sesión no exitoso.
\end{enumerate}

Si dichos tres flujos funcionan correctamente, entonces la práctica se puede considerar como correcta 

\section{Preguntas}

\begin{enumerate}
\item En el esquema de la base de datos proporcionada por los ayudantes, las contraseñas se guardan en texto plano, es decir, tal cual están almacenadas en la columna \texttt{password} corresponde a la contraseña de ese usuario, ¿por qué esto resulta inseguro desde el punto de vista de la ciberseguridad?
\item Explica cómo es que una aplicación front-end debería conectarse con tu aplicación de Flask para que los templates no sean necesarios. Supón que tu aplicación no utiliza la función \texttt{render\_template} ni \texttt{redirect}.
\item ¿Qué son los \texttt{CORS}?
\item En el contexto de aplicaciones back-end, ¿A qué nos referimos con \textit{manejo de secretos}?
\item Hablando del protocolo HTTP, describe un caso de uso en el que usar un método \texttt{GET} o un método \texttt{POST} no sea lo más ideal, es decir, un escenario en el que convenga usar un método que no sea ninguno de los mencionados.
\end{enumerate}

\section{Criterios de evaluación}

Para esta práctica se evaluarán de manera equitativa tanto los flujos implementados por el alumno, como las preguntas, es decir los $3$ flujos corresponderán al $50\%$ de la calificación y las $5$ preguntas al porcentaje restante. El caracter de cada reactivo será de criterio correcto o incorrecto, es decir, $0$ o $1$ o cualquier cosa que se le asimile.\\

\section{Lineamientos de entrega}

Todo deberá de estar en una carpeta comprimida (cualquier formato), incluyendo el código fuente, con el nombre: \texttt{Practica4\_<ApPat><ApMat>}, obviamente poniendo los respectivos apellidos paterno y materno.\\ 

La fecha de entrega para esta práctica es para el día 7 de Abril del 2026 antes de las 23:59:59.

\end{document}
